\section{AP-E11: compatible complete-minor cochain action}

AP-E11 replaces, rather than rescales, the AP-E10 stencil.  On an ordered
tetrahedral complex define the Alexander--Whitney front/back cup by
\begin{equation}
 (\alpha\smile_h\beta)[v_0\ldots v_{p+q}]
 =\alpha[v_0\ldots v_p]\,\beta[v_p\ldots v_{p+q}].
 \label{eq:e11-cup}
\end{equation}
The exact differential graded identities are
\begin{equation}
 d^2=0,\qquad d(\alpha\smile_h\beta)=d\alpha\smile_h\beta
 +(-1)^p\alpha\smile_h d\beta,
 \qquad (\alpha\smile_h\beta)\smile_h\gamma
 =\alpha\smile_h(\beta\smile_h\gamma).
 \label{eq:e11-dga}
\end{equation}
For $a^A=dn^A$, the normalized antisymmetric cochains
\begin{align}
 X_h^{AB}&=\frac12(a^A\smile_h a^B-a^B\smile_h a^A),\label{eq:e11-x}\\
 Z_h^{ABC}&=\frac1{3!}\sum_{\pi\in S_3}\operatorname{sgn}(\pi)
 a^{\pi(A)}\smile_h a^{\pi(B)}\smile_h a^{\pi(C)}
 \label{eq:e11-z}
\end{align}
are exactly the de Rham cochains of the P1 affine two- and three-minors.
The factor $1/2$ in \eqref{eq:e11-x} is the reference-triangle area.

For a $k$-simplex $\sigma=[i_0\ldots i_k]$, use the Whitney form
\begin{equation}
 \mathcal W_h(\sigma)=k!\sum_{r=0}^k(-1)^r\lambda_{i_r}
 d\lambda_{i_0}\wedge\cdots\widehat{d\lambda_{i_r}}\cdots
 \wedge d\lambda_{i_k}
 \label{eq:e11-whitney}
\end{equation}
and the consistent Hodge Gram matrix
$(M_k^T)_{\sigma\tau}=\int_T\mathcal W_h(\sigma)\wedge
\star\mathcal W_h(\tau)$.  Every $M_k^T$, $k=1,2,3$, is positive definite;
rigid rotations preserve it; and $\mathcal W_hd=d\mathcal W_h$ reconstructs
affine exterior forms exactly
\cite{Whitney1957,Dodziuk1976,ArnoldFalkWinther2010}.  The audit gives
$\lambda_{\min}=2.485\times10^{-2}$, rotation residual
$3.55\times10^{-15}$, and affine reconstruction residual
$2.40\times10^{-16}$.

The same third-cup family supplies topology.  On $T=[0123]$,
\begin{equation}
 C_h[n](T)=\epsilon_{ABCD}
 (n^A\smile_h a^B\smile_h a^C\smile_h a^D)(T)
 =\det[n_0\ n_1\ n_2\ n_3].
 \label{eq:e11-cup-det}
\end{equation}
If every tetrahedral vertex pair obeys $n_i\cdot n_j>\epsilon>-1/3$, then
for $v_T=\sum_i\lambda_i n_i$,
\begin{equation}
 |v_T|^2\ge\epsilon+(1-\epsilon)\sum_i\lambda_i^2
 \ge\frac{1+3\epsilon}{4}>0.
 \label{eq:e11-gap}
\end{equation}
Consequently $u_T=v_T/|v_T|$ is defined and its exact topological cochain is
\begin{equation}
 \Omega_h(T)=\frac{s_T}{2\pi^2}C_h[n](T)
 \int_{\Delta^3}|v_T(\lambda)|^{-4}\,d\lambda,\qquad
 \sum_T\Omega_h(T)=\deg u_h.
 \label{eq:e11-degree}
\end{equation}
This is the pullback formula for the $S^3$ volume form; the numerator is
constant on each affine tetrahedron and equals \eqref{eq:e11-cup-det}.

The compatible action is
\begin{equation}
 \begin{split}
 E_h&=\sum_T\int_T\left\{
 \frac12|Dv_h|^2+\frac R2|\wedge^2Dv_h|^2
 +\frac K2|\wedge^3Dv_h|^2+m^2(1-v_h^0)\right\}d^3x,\\
 &(R,K,m^2)=(1,0.35,0.12).
 \end{split}
 \label{eq:e11-action}
\end{equation}
The positive $K$ term is the Hodge norm of the complete family
\eqref{eq:e11-z}; its contraction with $n$ and radial dressing give
\eqref{eq:e11-degree}.  It closes the finite-grid topology-collapse mode
without reintroducing a mesh-direction quartic.

Every periodic P1 corrector preserves the means of all first, second, and
third minors.  Convex Jensen inequality on the complete minor vector proves
the exact cell formula
\begin{equation}
 \boxed{Q_{R,K}^{\rm cell}(A)
 =\frac12|A|^2+\frac R2|M_2(A)|^2+\frac K2|M_3(A)|^2,\qquad
 \varphi_{\min}=0.}
 \label{eq:e11-cell}
\end{equation}
It is identical on the uniform six-tet and both checkerboard five-tet meshes;
the optimized residual is $2.56\times10^{-14}$ and the minor-mean residual is
$1.78\times10^{-15}$.  This is the model-specific polyconvex mechanism
\cite{Ball1977}, not a scalar rescaling of AP-E10.

The Dirichlet term gives $H^1$ equicoercivity and forces the P1 limit back to
$S^3$.  The $L^2$ Hodge bounds on the compatible second and third minors,
together with Piola/null-Lagrangian identities, identify their weak limits.
Equation \eqref{eq:e11-gap} then identifies the normalized current
$\det[v_h,D_1v_h,D_2v_h,D_3v_h]/|v_h|^4$ and closes degree.  Convex lower
semicontinuity gives the liminf; smooth nodal interpolation gives recovery.
Thus the discrete theory Gamma-converges unconditionally to the fixed-degree
lower-semicontinuous relaxation.  Equality with the naive unrelaxed graph-
norm functional still requires a topology-sensitive smooth-density theorem
\cite{Bethuel1991}.

The production card passes $11/11$.  All thirteen unanchored cases have
$B=(1,1,1)$, projected gradient density at most $1.54\times10^{-7}$, and pair
dot at least $0.2724$.  The tail reaches $(N,L,a)=(37,8,0.222222)$; translation
and mesh comparisons use $(29,7,0.25)$.  With the same-action positive energy
density as quotient measure, the diagnostics are
\begin{center}
\begin{tabular}{lrrr}
\toprule
diagnostic&observed&threshold&pass\\
\midrule
joint-tail energy spread&0.437\%&3\%&yes\\
joint-tail centered-radius spread&3.489\%&5\%&yes\\
joint-tail radial CDF distance&4.067\%&5\%&yes\\
translated-start energy spread&$8.33\times10^{-12}$&1\%&yes\\
translated-start centered-radius spread&$7.16\times10^{-7}$&2\%&yes\\
five-/six-tet energy spread&0.463\%&3\%&yes\\
five-/six-tet centered-radius spread&0.752\%&5\%&yes\\
five-/six-tet radial CDF distance&3.556\%&5\%&yes\\
\bottomrule
\end{tabular}
\end{center}
Therefore
\begin{equation}
 \boxed{\begin{gathered}
 C_{\rm cell}=C_{\rm quotient}
 =C_{\rm relaxed\ regulator}=\mathrm{true},\\
 C_{\rm graph\ density}=C_{\rm classical\ regulator}
 =C_{\rm Hessian}=C_{\det}=\mathrm{false}.
 \end{gathered}}
 \label{eq:e11-decision}
\end{equation}
The AP-E10 stencil and mesh-universality blockers are closed.  The next
single mathematical gate is either fixed-degree smooth density in the
complete-minor graph norm or independent regularity and isolation of the
selected relaxed minimizer.  The classical same-action Hessian and regulated
determinant remain embargoed until one of those routes succeeds.
